\section{結晶・分子構造}

\subsection{令和3年度 問22}
\label{q:r03-reg-applied-22}

\noindent\textbf{原典：} sources/R03/17\_applied\_science.pdf、PDF物理8頁。

\textbf{公式正答：} 2（single）。

\paragraph{問題の要旨} ダイヤモンドと石墨の関係に相当しない鉱物の組合せを選ぶ。

\Needspace{4\baselineskip}\paragraph{確認した条件・記述（独自要約）}\mbox{}\par

\noindent\textbf{関係：} 化学組成が同じで結晶構造が異なる多形。\par

\noindent\textbf{方向：} 最も不適切な組合せ。\par

\paragraph{問われる核心} 多形と固溶体の区別

\paragraph{必要知識} 組成と結晶構造、同素体、斜長石端成分、代表的多形。

\paragraph{考え方・正答を選ぶ理由（査読待ち初稿）} 多形では同じ化学組成に異なる原子配置が対応する。曹長石と灰長石は化学組成が異なる斜長石の端成分なので2を選ぶ。結晶構造の型を保ちつつ組成が変わる固溶体と、同じ組成で構造が変わる多形とを軸で分ける。

\Needspace{4\baselineskip}\paragraph{選択肢の検討（査読待ち初稿）}\mbox{}\par

\noindent\textbf{1：多形に該当} 紅柱石／珪線石。両者ともAl2SiO5で、結晶構造が異なる。\par

\noindent\textbf{2：多形に非該当} 曹長石／灰長石。NaAlSi3O8とCaAl2Si2O8という異なる端成分であり、斜長石固溶体を作る。\par

\noindent\textbf{3：多形に該当} 方解石／あられ石。両者ともCaCO3で、結晶構造が異なる。\par

\noindent\textbf{4：多形に該当} 石英／トリディマイト。両者ともSiO2で、結晶構造が異なる。\par

\noindent\textbf{5：多形に該当} 閃亜鉛鉱／ウルツ鉱。両者ともZnSで、結晶構造が異なる。\par

\paragraph{再利用する考え方} 観察モード・切断方位・組成を確認し、多形と固溶体の区別を他の標本にも適用する。

\paragraph{共通理論} \hyperref[theory:mineral:composition]{第\ref*{theory:mineral:composition}節}

\par 参考：\cite{hom_albite,hom_anorthite,hom_andalusite,hom_calcite,hom_quartz,hom_sphalerite}。

\par 難易度 2（校準初稿）、優先度未評価。

\section{鉱物・岩石}

\subsection{平成24年度 問24}
\label{q:h24-reg-applied-24}

\noindent\textbf{原典：} sources/H24/17\_applied\_science.pdf、PDF物理8頁。

\textbf{公式正答：} 1（single）。

\paragraph{問題の要旨} 偏光顕微鏡を用いる同定が最も難しい鉱物を選ぶ。

\Needspace{4\baselineskip}\paragraph{確認した条件・記述（独自要約）}\mbox{}\par

\noindent\textbf{方向：} 偏光顕微鏡による同定が最も難しいもの。\par

\noindent\textbf{観察法：} 原典に透過光の明記なし。通常の岩石薄片の透過観察として解釈する。\par

\paragraph{問われる核心} 不透明鉱物と透過・反射観察の使い分け

\paragraph{必要知識} 硫化物、岩石薄片、オープン・クロスニコル、反射顕微鏡。

\paragraph{考え方・正答を選ぶ理由（査読待ち初稿）} 通常の岩石薄片を透過光で見る観察を想定すると、不透明な黄鉄鉱が最も不利であるため1を選ぶ。他の候補は透過光学的特徴を利用できる。原典は透過光と明記しておらず、反射偏光顕微鏡なら黄鉄鉱も観察できるので、あらゆる偏光観察で同定不能という意味には拡張しない。

\Needspace{4\baselineskip}\paragraph{選択肢の検討（査読待ち初稿）}\mbox{}\par

\noindent\textbf{1：条件に一致} 黄鉄鉱。不透明な硫化鉱物である。通常の透過薄片では内部の光学性を調べにくく、反射光観察等を要する。\par

\noindent\textbf{2：条件に不一致} 普通角閃石。通常の薄片では透過光による多色性、劈開、消光などを観察できる。\par

\noindent\textbf{3：条件に不一致} 単斜輝石。通常の薄片では屈折率、劈開、消光などが同定の手掛かりとなる。\par

\noindent\textbf{4：条件に不一致} 黒雲母。通常の薄片で強い多色性と一方向の劈開を観察できる。\par

\noindent\textbf{5：条件に不一致} ジルコン。小粒でも透明な薄片で高い屈折率・複屈折等を観察できる。\par

\paragraph{再利用する考え方} 観察モード・切断方位・組成を確認し、不透明鉱物と透過・反射観察の使い分けを他の標本にも適用する。

\paragraph{共通理論} \hyperref[theory:mineral:observation]{第\ref*{theory:mineral:observation}節}

\par 参考：\cite{hom_pyrite}。

\par 難易度 2（校準初稿）、優先度未評価。

\subsection{平成25年度 問22}
\label{q:h25-reg-applied-22}

\noindent\textbf{原典：} sources/H25/17\_applied\_science.pdf、PDF物理9頁。

\textbf{公式正答：} 3（single）。

\par\textbf{検討未確定：} 原典は黒雲母の組成又は代表密度を指定しない。MSA Handbookの黒雲母密度2.7〜3.3g/cm\textsuperscript{3}は設問境界2.85をまたぐため、1を一律に重鉱物とする根拠が不足する。独立選択は最も明確に境界未満の石英3。選択肢の単一性に関する疑義を残し部分校正とする。

\paragraph{問題の要旨} 比重2.85以上という重鉱物の定義を使い、重鉱物に分類されないものを選ぶ。

\Needspace{4\baselineskip}\paragraph{確認した条件・記述（独自要約）}\mbox{}\par

\noindent\textbf{境界：} 原典は重鉱物を比重2.85以上と定義する。\par

\noindent\textbf{方向：} 重鉱物に分類されないもの。\par

\paragraph{問われる核心} 密度による重鉱物区分と組成依存性

\paragraph{必要知識} 比重と密度の対応、分離用密度境界、固溶体の物性幅。

\paragraph{考え方・正答を選ぶ理由（査読待ち初稿）} 石英の代表比重は約2.65で、2.85以上という定義を満たさないので3を独立に選ぶ。ただし黒雲母は組成依存で2.7〜3.3という範囲があり、1も無条件に重鉱物とは断定できない。最も明確な軽鉱物と、問題文が省いた組成条件とを分けて記録する。

\Needspace{4\baselineskip}\paragraph{選択肢の検討（査読待ち初稿）}\mbox{}\par

\noindent\textbf{1：組成次第で境界をまたぐ} 黒雲母。黒雲母の密度は2.7〜3.3g/cm\textsuperscript{3}と報告され、2.85を下回る組成もある。一律の分類には留保が必要である。\par

\noindent\textbf{2：重鉱物に該当} かんらん石。Mg端成分のフォルステライトでも密度約3.275g/cm\textsuperscript{3}であり、境界より高い。\par

\noindent\textbf{3：明確に軽鉱物} 石英。石英の代表密度約2.65g/cm\textsuperscript{3}は境界より低い。\par

\noindent\textbf{4：重鉱物に該当} 普通輝石。普通輝石の報告密度3.19〜3.56g/cm\textsuperscript{3}は境界より高い。\par

\noindent\textbf{5：重鉱物に該当} 磁鉄鉱。磁鉄鉱の代表密度約5.175g/cm\textsuperscript{3}は境界より高い。\par

\paragraph{再利用する考え方} 観察モード・切断方位・組成を確認し、密度による重鉱物区分と組成依存性を他の標本にも適用する。

\paragraph{共通理論} \hyperref[theory:mineral:density]{第\ref*{theory:mineral:density}節}, \hyperref[theory:mineral:composition]{第\ref*{theory:mineral:composition}節}

\par 参考：\cite{hom_biotite,hom_quartz,hom_forsterite,hom_augite,hom_magnetite}。

\par 難易度 2（校準初稿）、優先度未評価。

\subsection{平成25年度 問23}
\label{q:h25-reg-applied-23}

\noindent\textbf{原典：} sources/H25/17\_applied\_science.pdf、PDF物理9頁。

\textbf{公式正答：} 4（single）。

\paragraph{問題の要旨} 造岩鉱物の偏光顕微鏡下の特徴から、最も不適切なものを選ぶ。

\Needspace{4\baselineskip}\paragraph{確認した条件・記述（独自要約）}\mbox{}\par

\noindent\textbf{方向：} 最も不適切。\par

\noindent\textbf{対象：} 通常の透過岩石薄片。消光角は適切な結晶切断面で判定する。\par

\paragraph{問われる核心} 劈開・屈折率・消光を組み合わせる同定

\paragraph{必要知識} 層状構造、オープンでの吸収色、クロスでの消光、組成累帯。

\paragraph{考え方・正答を選ぶ理由（査読待ち初稿）} 黒雲母は層状ケイ酸塩で、層に平行な一方向の劈開が明瞭である。二方向という4が誤りとなる。他の肢は通常の薄片観察の特徴に合う。直消光・斜消光の識別では切断方位により見掛け角が変わることを忘れない。

\Needspace{4\baselineskip}\paragraph{選択肢の検討（査読待ち初稿）}\mbox{}\par

\noindent\textbf{1：通常の薄片で正しい} 石英は無色で多色性がない。通常の石英薄片は無色であり、方向に依存する吸収色の変化は目立たない。\par

\noindent\textbf{2：正しい} カリ長石の屈折率は石英より低い。正長石の主屈折率約1.518〜1.525は石英の1.544〜1.553より低い。\par

\noindent\textbf{3：正しい} 斜長石には累帯構造がしばしば見られる。成長中の組成変化が結晶の内外の光学性に現れる。\par

\noindent\textbf{4：誤り} 黒雲母には2方向の明瞭な劈開がある。黒雲母の卓越する劈開は底面(001)の一方向である。\par

\noindent\textbf{5：通常の特徴として正しい／切断面依存} 普通角閃石は直消光しない。普通角閃石は通常斜消光を特徴とするが、見掛けの角度は切断方位に依存するので任意の一粒だけで断定しない。\par

\paragraph{再利用する考え方} 観察モード・切断方位・組成を確認し、劈開・屈折率・消光を組み合わせる同定を他の標本にも適用する。

\paragraph{共通理論} \hyperref[theory:mineral:cleavage]{第\ref*{theory:mineral:cleavage}節}, \hyperref[theory:mineral:observation]{第\ref*{theory:mineral:observation}節}, \hyperref[theory:mineral:extinction]{第\ref*{theory:mineral:extinction}節}, \hyperref[theory:mineral:texture]{第\ref*{theory:mineral:texture}節}

\par 参考：\cite{hom_biotite,hom_quartz,hom_orthoclase,nelson_inosilicates}。

\par 難易度 2（校準初稿）、優先度未評価。

\subsection{平成26年度 問28}
\label{q:h26-reg-applied-28}

\noindent\textbf{原典：} sources/H26/17\_applied\_science.pdf、PDF物理12頁。

\textbf{公式正答：} 4（single）。

\paragraph{問題の要旨} 花こう岩を構成する鉱物として最も一般的でないものを選ぶ。

\Needspace{4\baselineskip}\paragraph{確認した条件・記述（独自要約）}\mbox{}\par

\noindent\textbf{方向：} 条件に最も該当する鉱物を選ぶ。\par

\paragraph{問われる核心} 珪長質岩と苦鉄質鉱物の組合せ

\paragraph{必要知識} 花こう岩、石英・長石、かんらん石、シリカ飽和度。

\paragraph{考え方・正答を選ぶ理由（査読待ち初稿）} 通常の花こう岩の構成は石英、カリ長石、斜長石と黒雲母などである。かんらん石は苦鉄質側を特徴付けるため4を選ぶ。稀な岩相や捕獲結晶の存在まで否定する問題ではなく、典型的な鉱物組合せを問う。

\Needspace{4\baselineskip}\paragraph{選択肢の検討（査読待ち初稿）}\mbox{}\par

\noindent\textbf{1：条件に不一致} 黒雲母。花こう岩に含まれる代表的な有色鉱物である。\par

\noindent\textbf{2：条件に不一致} 石英。花こう岩の主要な構成鉱物である。\par

\noindent\textbf{3：条件に不一致} カリ長石。花こう岩の主要な構成鉱物である。\par

\noindent\textbf{4：条件に一致} かんらん石。かんらん石は通常、苦鉄質・超苦鉄質岩を特徴付け、石英に富む一般的な花こう岩の主成分ではない。\par

\noindent\textbf{5：条件に不一致} 斜長石。花こう岩の主要な構成鉱物である。\par

\paragraph{再利用する考え方} 観察モード・切断方位・組成を確認し、珪長質岩と苦鉄質鉱物の組合せを他の標本にも適用する。

\paragraph{共通理論} \hyperref[theory:mineral:assemblage]{第\ref*{theory:mineral:assemblage}節}

\par 参考：\cite{usgs_igneous_components,hom_forsterite}。

\par 難易度 2（校準初稿）、優先度未評価。

\subsection{平成27年度 問30}
\label{q:h27-reg-applied-30}

\noindent\textbf{原典：} sources/H27/17\_applied\_science.pdf、PDF物理11頁。

\textbf{公式正答：} 2（single）。

\paragraph{問題の要旨} 流紋岩・安山岩・玄武岩に共通する主要な造岩鉱物を選ぶ。

\Needspace{4\baselineskip}\paragraph{確認した条件・記述（独自要約）}\mbox{}\par

\noindent\textbf{方向：} 条件に最も該当する鉱物を選ぶ。\par

\paragraph{問われる核心} 斜長石固溶体と火成岩系列

\paragraph{必要知識} 流紋岩・安山岩・玄武岩、斜長石のNa-Ca組成、典型的鉱物組合せ。

\paragraph{考え方・正答を選ぶ理由（査読待ち初稿）} 岩石系列をまたいで残るのは斜長石2である。アルバイト成分に富む側からアノーサイト成分に富む側へ組成が変わるため、岩石名が変わっても斜長石という鉱物群は広く現れる。個々の急冷ガラス質・無斑晶試料に必ず可視結晶があるという全称命題にはしない。

\Needspace{4\baselineskip}\paragraph{選択肢の検討（査読待ち初稿）}\mbox{}\par

\noindent\textbf{1：条件に不一致} 石英。石英は流紋岩などで一般的だが、通常の玄武岩に共通する主要鉱物ではない。\par

\noindent\textbf{2：条件に一致} 斜長石。Na側からCa側まで組成が変わる斜長石は、珪長質から苦鉄質まで広く出現する。\par

\noindent\textbf{3：条件に不一致} 黒雲母。黒雲母は一般的な玄武岩に共通する主要鉱物ではない。\par

\noindent\textbf{4：条件に不一致} 輝石。輝石は苦鉄質〜中性岩で一般的だが、通常の流紋岩に共通する主要鉱物とはいえない。\par

\noindent\textbf{5：条件に不一致} かんらん石。かんらん石は苦鉄質岩側で特徴的で、流紋岩には一般的でない。\par

\paragraph{再利用する考え方} 観察モード・切断方位・組成を確認し、斜長石固溶体と火成岩系列を他の標本にも適用する。

\paragraph{共通理論} \hyperref[theory:mineral:composition]{第\ref*{theory:mineral:composition}節}, \hyperref[theory:mineral:assemblage]{第\ref*{theory:mineral:assemblage}節}

\par 参考：\cite{hom_albite,hom_anorthite,usgs_igneous_components}。

\par 難易度 2（校準初稿）、優先度未評価。

\subsection{平成29年度 問27}
\label{q:h29-reg-applied-27}

\noindent\textbf{原典：} sources/H29/17\_applied\_science.pdf、PDF物理9頁。

\textbf{公式正答：} 4（single）。

\paragraph{問題の要旨} 花こう岩を構成する鉱物として最も一般的でないものを選ぶ。

\Needspace{4\baselineskip}\paragraph{確認した条件・記述（独自要約）}\mbox{}\par

\noindent\textbf{方向：} 条件に最も該当する鉱物を選ぶ。\par

\paragraph{問われる核心} 珪長質岩と苦鉄質鉱物の組合せ

\paragraph{必要知識} 花こう岩、石英・長石、かんらん石、シリカ飽和度。

\paragraph{考え方・正答を選ぶ理由（査読待ち初稿）} 通常の花こう岩の構成は石英、カリ長石、斜長石と黒雲母などである。かんらん石は苦鉄質側を特徴付けるため4を選ぶ。稀な岩相や捕獲結晶の存在まで否定する問題ではなく、典型的な鉱物組合せを問う。

\Needspace{4\baselineskip}\paragraph{選択肢の検討（査読待ち初稿）}\mbox{}\par

\noindent\textbf{1：条件に不一致} 黒雲母。花こう岩に含まれる代表的な有色鉱物である。\par

\noindent\textbf{2：条件に不一致} 石英。花こう岩の主要な構成鉱物である。\par

\noindent\textbf{3：条件に不一致} カリ長石。花こう岩の主要な構成鉱物である。\par

\noindent\textbf{4：条件に一致} かんらん石。かんらん石は通常、苦鉄質・超苦鉄質岩を特徴付け、石英に富む一般的な花こう岩の主成分ではない。\par

\noindent\textbf{5：条件に不一致} 斜長石。花こう岩の主要な構成鉱物である。\par

\paragraph{再利用する考え方} 観察モード・切断方位・組成を確認し、珪長質岩と苦鉄質鉱物の組合せを他の標本にも適用する。

\paragraph{共通理論} \hyperref[theory:mineral:assemblage]{第\ref*{theory:mineral:assemblage}節}

\par 参考：\cite{usgs_igneous_components,hom_forsterite}。

\par 難易度 2（校準初稿）、優先度未評価。

\subsection{平成30年度 問31}
\label{q:h30-reg-applied-31}

\noindent\textbf{原典：} sources/H30/17\_applied\_science.pdf、PDF物理10頁。

\textbf{公式正答：} 4（single）。

\paragraph{問題の要旨} 造岩鉱物の薄片光学特性から最も不適切なものを選ぶ。

\Needspace{4\baselineskip}\paragraph{確認した条件・記述（独自要約）}\mbox{}\par

\noindent\textbf{方向：} 最も不適切。\par

\noindent\textbf{観察：} 通常の切断面の透過偏光観察。\par

\paragraph{問われる核心} 晶系と消光方向の関係

\paragraph{必要知識} 斜方晶系・単斜晶系、切断方位、角閃石の劈開、双晶・波動消光。

\paragraph{考え方・正答を選ぶ理由（査読待ち初稿）} 輝石の晶系と消光の対応が逆になっている4を選ぶ。斜方輝石の直消光、単斜輝石の斜消光を基本にするが、単斜輝石の特定の切断面でも直消光があり得る。劈開角、屈折率、双晶等の情報を合わせて判断する。

\Needspace{4\baselineskip}\paragraph{選択肢の検討（査読待ち初稿）}\mbox{}\par

\noindent\textbf{1：正しい} 石英では波動消光が見られることがある。結晶内部の歪みで光学方位が連続的に変わると一粒の消光が一斉でなくなる。\par

\noindent\textbf{2：正しい} 普通角閃石には2方向の劈開が見られる。角閃石の柱面劈開は約56°・124°で交わる二方向である。\par

\noindent\textbf{3：通常の特徴として正しい} かんらん石は直消光し劈開に乏しい。通常のかんらん石は直消光で、劈開は弱く顕著ではない。\par

\noindent\textbf{4：誤り} 斜方輝石は斜消光し単斜輝石は直消光する。通常の特徴は斜方輝石が直消光、単斜輝石が斜消光であり、記述が逆である。単斜輝石も切断面次第で平行消光を示す。\par

\noindent\textbf{5：正しい} 斜長石にはアルバイト双晶や累帯構造が見られる。斜長石は反復双晶と成長中の組成変化を示し得る。\par

\paragraph{再利用する考え方} 観察モード・切断方位・組成を確認し、晶系と消光方向の関係を他の標本にも適用する。

\paragraph{共通理論} \hyperref[theory:mineral:cleavage]{第\ref*{theory:mineral:cleavage}節}, \hyperref[theory:mineral:extinction]{第\ref*{theory:mineral:extinction}節}, \hyperref[theory:mineral:texture]{第\ref*{theory:mineral:texture}節}

\par 参考：\cite{nelson_inosilicates,nelson_biaxial,hom_albite}。

\par 難易度 2（校準初稿）、優先度未評価。

\subsection{令和元年度 問26}
\label{q:r01-reg-applied-26}

\noindent\textbf{原典：} sources/R01/17\_applied\_science.pdf、PDF物理9頁。

\textbf{公式正答：} 4（single）。

\paragraph{問題の要旨} 偏光顕微鏡で使う用語の説明から最も不適切なものを選ぶ。

\Needspace{4\baselineskip}\paragraph{確認した条件・記述（独自要約）}\mbox{}\par

\noindent\textbf{方向：} 最も不適切。\par

\noindent\textbf{区別：} オープンニコルは検光子を外し、クロスニコルは偏光子と検光子を直交させる。\par

\paragraph{問われる核心} 吸収異方性と複屈折による干渉の区別

\paragraph{必要知識} 偏光子・検光子、ベッケ線、複屈折・干渉色、消光、双晶。

\paragraph{考え方・正答を選ぶ理由（査読待ち初稿）} 多色性は吸収の異方性によってオープンニコルで現れる。クロスニコルでの干渉色の変化を多色性と呼ぶ4が誤りである。波動としての位相遅れと、方向ごとの吸収係数の違いとを分ければ用語を整理できる。

\Needspace{4\baselineskip}\paragraph{選択肢の検討（査読待ち初稿）}\mbox{}\par

\noindent\textbf{1：正しい} オープンニコルのベッケ線は隣接する鉱物の屈折率比較に使う。焦点をずらしたときの明るい境界線の移動を利用して屈折率の高低を比べる。\par

\noindent\textbf{2：正しい} クロスニコルの干渉色は複屈折で生じる。異なる固有偏光の光路差が、検光子で合成される際の波長依存強度を生む。\par

\noindent\textbf{3：正しい} クロスニコルで鉱物が暗くなる状態を消光という。透過強度が小さくなり暗く見える状態である。\par

\noindent\textbf{4：誤り} クロスニコルで回転により鉱物の色が変わる性質を多色性という。多色性はオープンニコルで偏光方向に応じ吸収色が変わる現象である。クロスでの干渉色変化と混同している。\par

\noindent\textbf{5：観察方位によって正しい} 双晶の結晶部分はクロスニコルで2つ以上の異なる色を示し得る。双晶ドメインは一定の関係で異なる方位を持ち、同じ回転角でも消光や色が異なり得る。常に全方位で色が異なるとは限らない。\par

\paragraph{再利用する考え方} 観察モード・切断方位・組成を確認し、吸収異方性と複屈折による干渉の区別を他の標本にも適用する。

\paragraph{共通理論} \hyperref[theory:mineral:observation]{第\ref*{theory:mineral:observation}節}, \hyperref[theory:mineral:extinction]{第\ref*{theory:mineral:extinction}節}, \hyperref[theory:mineral:texture]{第\ref*{theory:mineral:texture}節}

\par 参考：\cite{nelson_biaxial}。

\par 難易度 2（校準初稿）、優先度未評価。

\subsection{令和元年度再試験 問25}
\label{q:r01-ret-applied-25}

\noindent\textbf{原典：} sources/R01\_retest/17\_applied\_science.pdf、PDF物理8頁。

\textbf{公式正答：} 3（single）。

\paragraph{問題の要旨} 偏光顕微鏡を用いる同定が最も難しい鉱物を選ぶ。

\Needspace{4\baselineskip}\paragraph{確認した条件・記述（独自要約）}\mbox{}\par

\noindent\textbf{方向：} 偏光顕微鏡による同定が最も難しいもの。\par

\noindent\textbf{観察法：} 原典に透過光の明記なし。通常の岩石薄片の透過観察として解釈する。\par

\paragraph{問われる核心} 不透明鉱物と透過・反射観察の使い分け

\paragraph{必要知識} 硫化物、岩石薄片、オープン・クロスニコル、反射顕微鏡。

\paragraph{考え方・正答を選ぶ理由（査読待ち初稿）} 通常の岩石薄片を透過光で見る観察を想定すると、不透明な黄鉄鉱が最も不利であるため3を選ぶ。他の候補は透過光学的特徴を利用できる。原典は透過光と明記しておらず、反射偏光顕微鏡なら黄鉄鉱も観察できるので、あらゆる偏光観察で同定不能という意味には拡張しない。

\Needspace{4\baselineskip}\paragraph{選択肢の検討（査読待ち初稿）}\mbox{}\par

\noindent\textbf{1：条件に不一致} 普通角閃石。通常の薄片では透過光による多色性、劈開、消光などを観察できる。\par

\noindent\textbf{2：条件に不一致} 斜長石。通常の薄片で双晶・消光・累帯構造等を観察できる。\par

\noindent\textbf{3：条件に一致} 黄鉄鉱。不透明な硫化鉱物である。通常の透過薄片では内部の光学性を調べにくく、反射光観察等を要する。\par

\noindent\textbf{4：条件に不一致} 単斜輝石。通常の薄片では屈折率、劈開、消光などが同定の手掛かりとなる。\par

\noindent\textbf{5：条件に不一致} 黒雲母。通常の薄片で強い多色性と一方向の劈開を観察できる。\par

\paragraph{再利用する考え方} 観察モード・切断方位・組成を確認し、不透明鉱物と透過・反射観察の使い分けを他の標本にも適用する。

\paragraph{共通理論} \hyperref[theory:mineral:observation]{第\ref*{theory:mineral:observation}節}

\par 参考：\cite{hom_pyrite}。

\par 難易度 2（校準初稿）、優先度未評価。

\subsection{令和3年度 問25}
\label{q:r03-reg-applied-25}

\noindent\textbf{原典：} sources/R03/17\_applied\_science.pdf、PDF物理9頁。

\textbf{公式正答：} 4（single）。

\paragraph{問題の要旨} 偏光顕微鏡による造岩鉱物の識別の記述から、最も不適切なものを選ぶ。

\Needspace{4\baselineskip}\paragraph{確認した条件・記述（独自要約）}\mbox{}\par

\noindent\textbf{方向：} 最も不適切。\par

\noindent\textbf{観察：} 通常の薄片・適切な切断面。\par

\paragraph{問われる核心} 劈開と結晶内歪みによる消光の区別

\paragraph{必要知識} 石英、黒雲母、輝石、角閃石、かんらん石の薄片特徴。

\paragraph{考え方・正答を選ぶ理由（査読待ち初稿）} 石英には普通の薄片で顕著な劈開がないので4が不適切である。波動消光は結晶内の歪みを反映するため石英にも生じる。複合文の一部だけが正しくても、肢全体を正しいと扱わない。

\Needspace{4\baselineskip}\paragraph{選択肢の検討（査読待ち初稿）}\mbox{}\par

\noindent\textbf{1：通常の特徴として正しい} 黒雲母は強い多色性を示し劈開線に平行な直消光をする。黒雲母は吸収異方性が強く、底面劈開線を消光方向の目安にできる。晶系の厳密な光学方位では小角のずれがあり得る。\par

\noindent\textbf{2：通常の特徴として正しい} 斜方輝石は直消光、普通輝石は斜消光する。識別の基本特徴に合う。切断面による見掛け角の違いは確認が必要である。\par

\noindent\textbf{3：通常の特徴として正しい} 普通角閃石は高屈折率で強い多色性を示す。普通角閃石の屈折率・多色性を併用する特徴に合う。\par

\noindent\textbf{4：前半が誤り} 石英は顕著な劈開を持ち、波動消光することがある。石英の劈開は通常顕著でない。波動消光し得るという後半だけは正しい。\par

\noindent\textbf{5：通常の特徴として正しい} かんらん石は直消光し劈開に乏しい。かんらん石の通常の特徴に合う。劈開が絶対に存在しないという意味ではない。\par

\paragraph{再利用する考え方} 観察モード・切断方位・組成を確認し、劈開と結晶内歪みによる消光の区別を他の標本にも適用する。

\paragraph{共通理論} \hyperref[theory:mineral:cleavage]{第\ref*{theory:mineral:cleavage}節}, \hyperref[theory:mineral:extinction]{第\ref*{theory:mineral:extinction}節}, \hyperref[theory:mineral:texture]{第\ref*{theory:mineral:texture}節}

\par 参考：\cite{hom_quartz,hom_biotite,hom_forsterite,nelson_inosilicates,nelson_biaxial}。

\par 難易度 2（校準初稿）、優先度未評価。

\subsection{令和4年度 問23}
\label{q:r04-reg-applied-23}

\noindent\textbf{原典：} sources/R04/17\_applied\_science.pdf、PDF物理8頁。

\textbf{公式正答：} 4（single）。

\paragraph{問題の要旨} 火成岩を作る鉱物の構造・光学性・分化の記述から最も不適切なものを選ぶ。

\Needspace{4\baselineskip}\paragraph{確認した条件・記述（独自要約）}\mbox{}\par

\noindent\textbf{方向：} 最も不適切。\par

\noindent\textbf{顕微鏡：} 透過薄片のオープン／クロスを区別する。\par

\paragraph{問われる核心} 結晶・非晶質の光学応答とマグマ分化

\paragraph{必要知識} SiO4四面体、層状構造、吸収異方性、等方性、結晶分離。

\paragraph{考え方・正答を選ぶ理由（査読待ち初稿）} 通常の火山ガラスは等方的な非晶質であり、歪んだ結晶のような波動消光を示さない。4を選ぶ。輝石の多色性の強さは種・組成に依存するので「全ての輝石が強い多色性を持つ」と暗記しない。結晶分化で変わるのは分離後の残液で、結晶を含めた閉鎖系全体の元素保存が破れるわけではない。

\Needspace{4\baselineskip}\paragraph{選択肢の検討（査読待ち初稿）}\mbox{}\par

\noindent\textbf{1：主要鉱物について正しい} 火成岩の主要鉱物はSiO4四面体を骨格とするケイ酸塩である。石英・長石・輝石・角閃石・かんらん石・雲母はケイ酸塩の主要な造岩鉱物である。酸化物等の副成分まで全てケイ酸塩とはしていない。\par

\noindent\textbf{2：正しい} 黒雲母の劈開は一方向である。層に平行な底面劈開が卓越する。\par

\noindent\textbf{3：通常の特徴として正しい／鉱物種・組成依存} オープンでは黒雲母・角閃石・輝石に多色性が見られ、かんらん石・ざくろ石には見られない。吸収異方性を持つ有色鉱物には多色性があり得るが、輝石の種・組成で弱い又は目立たない場合もある。ざくろ石は通常等方性、かんらん石は通常薄片で無色である。\par

\noindent\textbf{4：誤り} 火山ガラスはクロスニコルで波動消光を示す。通常の火山ガラスは光学的等方性で、クロスでは回転しても暗い。歪んだ石英の波動消光とは違う。\par

\noindent\textbf{5：正しい} 晶出した鉱物を分離する結晶分化でマグマ組成が変わる。晶出相を残液から除くと、除去された元素比に応じ残液組成が変わる。\par

\paragraph{再利用する考え方} 観察モード・切断方位・組成を確認し、結晶・非晶質の光学応答とマグマ分化を他の標本にも適用する。

\paragraph{共通理論} \hyperref[theory:mineral:structure]{第\ref*{theory:mineral:structure}節}, \hyperref[theory:mineral:observation]{第\ref*{theory:mineral:observation}節}, \hyperref[theory:mineral:extinction]{第\ref*{theory:mineral:extinction}節}, \hyperref[theory:mineral:assemblage]{第\ref*{theory:mineral:assemblage}節}

\par 参考：\cite{hom_biotite,hom_augite,nelson_inosilicates,usgs_igneous_components}。

\par 難易度 2（校準初稿）、優先度未評価。

\subsection{令和6年度 問23}
\label{q:r06-reg-applied-23}

\noindent\textbf{原典：} sources/R06/17\_applied\_science.pdf、PDF物理8頁。

\textbf{公式正答：} 2（single）。

\paragraph{問題の要旨} 火成岩の造岩鉱物と結晶構造に関する5記述のうち、最も不適切なものを選ぶ。

\Needspace{4\baselineskip}\paragraph{確認した条件・記述（独自要約）}\mbox{}\par

\noindent\textbf{方向：} 最も不適切。\par

\noindent\textbf{主語：} 原典は火成岩。OCRの別語へ置き換えない。\par

\paragraph{問われる核心} 構造と組成の二軸による鉱物の区別

\paragraph{必要知識} ケイ酸塩構造、多形・固溶体、底面劈開、結晶化系列、密度。

\paragraph{考え方・正答を選ぶ理由（査読待ち初稿）} 2は固溶体の定義を多形と取り違えている。多形は同じ化学組成で構造が異なり、固溶体は構造の型を保って元素置換により組成が変わる。黒雲母の劈開は層状構造から、密度比較は組成と原子の詰まり方から理解する。

\Needspace{4\baselineskip}\paragraph{選択肢の検討（査読待ち初稿）}\mbox{}\par

\noindent\textbf{1：主要鉱物について正しい} 火成岩の主要な造岩鉱物はSiO4四面体を骨格とするケイ酸塩である。主要なケイ酸塩造岩鉱物を述べている。副成分を含む全鉱物という主張ではない。\par

\noindent\textbf{2：誤り} 結晶構造が同じで組成が連続的に変わる鉱物を多形という。構造の型を保ち組成が変わるものは固溶体である。多形は同じ組成に異なる結晶構造が対応する。\par

\noindent\textbf{3：正しい} 黒雲母はSiO4四面体の層状構造を持ち、層間の結合が弱いため劈開が明瞭である。四面体層を含む複合層と層間の結合の差が、底面劈開を生む。\par

\noindent\textbf{4：通常の結晶化系列で正しい} 玄武岩質マグマが冷えると、有色鉱物の中では融点が高いかんらん石が最初に晶出する。通常のボーエン系列でかんらん石は高温側の初期鉱物である。実際の晶出は圧力・含水量・組成にも依存する。\par

\noindent\textbf{5：代表的物性で正しい} かんらん石・普通輝石・普通角閃石・黒雲母・石英の中では石英の比重が最も小さい。石英は約2.65で、候補の普通の有色鉱物より小さい。黒雲母の密度幅にも留意するが下端約2.7は石英より大きい。\par

\paragraph{再利用する考え方} 観察モード・切断方位・組成を確認し、構造と組成の二軸による鉱物の区別を他の標本にも適用する。

\paragraph{共通理論} \hyperref[theory:mineral:structure]{第\ref*{theory:mineral:structure}節}, \hyperref[theory:mineral:composition]{第\ref*{theory:mineral:composition}節}, \hyperref[theory:mineral:cleavage]{第\ref*{theory:mineral:cleavage}節}, \hyperref[theory:mineral:assemblage]{第\ref*{theory:mineral:assemblage}節}, \hyperref[theory:mineral:density]{第\ref*{theory:mineral:density}節}

\par 参考：\cite{hom_quartz,hom_biotite,hom_forsterite,hom_augite}。

\par 難易度 2（校準初稿）、優先度未評価。

\section{未確定}

\subsection{平成25年度 問18}
\label{q:h25-reg-applied-18}

\noindent\textbf{原典：} sources/H25/17\_applied\_science.pdf、PDF物理8頁。

\textbf{公式正答：} 3（single）。

\par\textit{原典本文の校正と独自解説は未作成。分類のみの登録項目。}
